\scanpage{22}
\deleted{On procède par récurrence sur $\operatorname{card} K$,} \inserted{On note d'abord que} \deleted{pour prouver que} (si $V=\underset{\substack{\rotatebox{90}{$=$}\\ \inserted{\bigcap_{\alpha\in K}\Sigma_\alpha^-}}}{\underbrace{X\setminus\bigcup_{\alpha\in K}\widehat{\Sigma}_\alpha^+}}$, les $\widehat{\Sigma}_\alpha^+$ \unsure{(}$\alpha\in K$ étant \unsure{disjoints|disjointes})
\[
\left\{\begin{aligned}
&\overline V\cap E_\alpha\ne\varnothing\\
&\makebox[0pt][r]{$\longrightarrow$\hspace{.15em}}V\cap\Sigma_\alpha^-\deleted{\unsure{\ne|=}\varnothing}\ne\varnothing
\qquad\text{alors que } V\cap\Sigma_\alpha^+=\varnothing
\end{aligned}\right.
\]
% The marginal arrow crosses the brace and ends at V in its second row.
\reversemarginpar\margin{\raggedleft\raisebox{.8\baselineskip}{trivial}}\normalmarginpar

De façon précise, je dis que
\[
\begin{aligned}
\overline V\cap E_\alpha & = \overline{\Sigma_\alpha^-}\cap E_\alpha\qquad\text{(qui est donc }\ne\varnothing\text{)}\\[-2pt]
\underset{\inserted{\overline V\cap\widehat{\Sigma}_\alpha^+}}{\rotatebox{90}{$=$}}&
\hspace{1em}\underset{\inserted{\overline{\Sigma_\alpha^-}\cap\widehat{\Sigma}_\alpha^+}}{\rotatebox{90}{$\unsure{\deleted{=}|=|\Rightarrow}$}}
\end{aligned}
\]

\deleted{Comme $V\subset\Sigma_\alpha^-$, on a $V\cap E_\alpha\subset\overline{\Sigma_\alpha^-}\cap E_\alpha$, prouvons l'inclusion inverse}
\[
\deleted{\overline{\Sigma_\alpha^-}\cap E_\alpha\subset\overline V\cap E_\alpha}
\]
\deleted{Soit $x\in\overline{\Sigma_\alpha^-}\cap E_\alpha$, prouvons $x\in\overline V\cap E_\alpha$ i.e. $x\in\overline V$, soit $W$ un voisinage de $x$, montrons qu'il rencontre}

\deleted{Or} \deleted{\unsure{on}} \deleted{peut considérer cette} \deleted{\unsure{formule|relation}} \deleted{\inserted{égalité} comme égalité dans l'espace}
\[
\deleted{\unsure{Y|Z}=\widehat{\Sigma}_\alpha^-=\Sigma_\alpha^-\cup E_\alpha}
\]
\deleted{dans lequel on a deux parties fermées disjointes}

Or dans $X$ on a les deux parties \uline{fermées} \uline{disjointes}
\[
Y=\widehat{\Sigma}_\alpha^+,\qquad Z=\bigcup_{\beta\in K\setminus\{\alpha\}}\widehat{\Sigma}_\beta^+
\]
et on doit prouver
\[
\overline{(X\setminus(Y\cup Z))}\cap Y=\overline{(X\setminus Y)}\cap Y
\]
ce qui est immédiat
% archive sheet number: 21
